\par\bigskip\Needspace{6\baselineskip}
\begin{center}\textls[100]{b. Inhalt der Abschnitte I—IV der D. A.}\end{center}
\begin{center}\bfseries 3.\end{center}
\noindent\hspace*{\parindent}Wenden wir uns nun zum Inhalte der D. A. selbst. Ihre vier ersten Abschnitte geben die Elemente der heutigen Zahlentheorie, genauer desjenigen Teils derselben, den man multiplikative Zahlentheorie nennt, weil er wesentlich auf die Darstellbarkeit der Zahlen als Produkten von einfachsten Faktoren, den Primzahlen, gegründet und auf die Teilung der Zahlen durch einander, die Eigenschaften der Reste, welche dabei verbleiben, u. dgl. mehr gerichtet ist. Hier fand \textsc{Gauss} Vieles und Wichtiges von dem, was er selbständig erforscht hatte, schon von älteren Mathematikern, wie \textsc{Fermat}, \textsc{Euler}, \textsc{Lagrange} und \textsc{Legendre}, teils nur bemerkt, teils aber auch bewiesen. Er aber hat es Alles methodisch und in wissenschaftlichem Zusammenhange uns neu entwickelt und dargestellt.

Gleich die Einführung des Begriffes der Kongruenz, mit welcher \textls[100]{der erste Abschnitt} beginnt, eines Begriffs, der zwar auch sonst schon stillschweigend verwendet, aber bis dahin noch nicht formuliert und methodisch verwertet worden war, gilt eine Tat: zwei Zahlen \(a,\,b\) heissen einander kongruent nach einem Modul \(m\), in Zeichen:
\[
a\equiv b\ (\mathrm{mod.}\ m),
\]
wenn sie gleichen Rest lassen bei der Teilung durch \(m\). Und äusserst glücklich gewählt ist das Zeichen für solches Verhalten, weil es die weitgreifende Analogie, die sich alsbald zwischen Kongruenzen und Gleichungen herausstellt,\newpage
\SourceLeaf{16}\ArithmeticHead{9}{R}{INHALT DER ABSCHNITTE I—IV DER D. A.}
\noindent zur lichten Anschauung bringt. So steht der Forderung, eine algebraische \textls[100]{Gleichung} \(n^{\mathrm{ten}}\) Grades
\[
ax^n+a_1x^{n-1}+\cdots+a_n=0
\]
aufzulösen, falls die Koeffizienten ganzzahlig gedacht werden, die Aufgabe zur Seite, die Wurzeln der \textls[100]{Kongruenz}
\[
ax^n+a_1x^{n-1}+\cdots+a_n\equiv0\ (\mathrm{mod.}\ m)
\]
d. h. die inkongruenten ganzzahligen Werte von \(x\) zu finden, die ihr genügen. \textls[100]{Der zweite Abschnitt} gibt neben Fundamentalsätzen über Teilbarkeit der Zahlen die Lösung jener Aufgabe für die Kongruenzen ersten Grades, welche im Grunde mit der schon von \textsc{Euler} und \textsc{Lagrange} nach verschiedenen Methoden behandelten Auflösung der unbestimmten Gleichung
\[
ax+by=c
\]
übereinkommt. Während \textsc{Gauss} die analoge Aufgabe für mehrere lineare Kongruenzen mit mehr als einer Unbekannten nur flüchtig behandelt, zieht er aus der Auffindung einer Zahl, welche nach mehreren gegebenen Moduln gegebene Reste lässt, eine neue Bestimmung des bereits von \textsc{Euler}\footnote{Novi Comm. Acad. Petrop. \textbf{8} (1760/61) 1763, S. 74, Opera Omnia, Ser. I, vol. 2, S. 531.} ermittelten Ausdrucks der Funktion \(\varphi(m)\), welche die Anzahl der zu \(m\) primen Zahlen, welche \(\EqualAboveLess m\), bezeichnet. Der Abschnitt schliesst mit dem nicht wesentlich von dem \textsc{Lagrange}schen\begingroup\renewcommand{\thefootnote}{**)}\footnote{Mémoires de l’Acad. Berlin (1768) 1770 \textbf{24}, S. 192, Oeuvres II, S. 667.}\endgroup\ verschiedenen Beweise der Tatsache, dass eine Kongruenz in Bezug auf einen Primzahlmodul nicht mehr Wurzeln haben kann, als ihr Grad beträgt.

Etwas unvermittelt schiebt sich im art. 42 jener fundamentale \textsc{Gauss}sche Satz ein, dass, wenn eine ganze Funktion von \(x\) mit ganzzahligen Koeffizienten in ein Produkt zweier ganzer Funktionen mit rationalen Koeffizienten zerlegbar ist, diese Koeffizienten ebenfalls ganzzahlig sein müssen, ein Satz, welcher die Grundlage für die besonders von \textsc{Kronecker} ausgebildete Arithmetik der ganzen Funktionen bildet und auch in andern Disziplinen, u. a. für einen der wichtigsten Sätze über Ideale eines algebraischen Zahlkörpers, sich grundlegend gezeigt hat. In \textsc{Gauss}’ \textit{Tagebuch} (Nr. 69) wird die Auffindung dieses Satzes
\HistoricalFoot{X\textsubscript{2}.}{2}
\clearpage\SourceLeaf{17}\ArithmeticHead{10}{L}{BACHMANN, ÜBER GAUSS’ ZAHLENTHEORETISCHE ARBEITEN.}
\noindent verhältnismässig spät, erst auf den 23. Juli 1797 datiert; vermutlich den Studien zur Theorie der höheren Kongruenzen, welche der achte Abschnitt behandelt, entsprungen, ist er hier eingefügt.
